\begin{examplebox}
	\textbf{\textcolor{CExample}{例1（基础）}}
	\textbf{\textcolor{CExample}{题目：}}
	计算 $2^{\log_2 3}$ 的值。
	\textbf{\textcolor{CExample}{解题步骤：}}
	\begin{description}[style=nextline, leftmargin=2.8em, labelsep=0.5em, itemsep=0.45em]
		\item[\textbf{第一步（识别底数）：}]
			观察底数：指数底数为 $2$，对数底数也为 $2$，底数相同且符合条件。
			\par\textbf{理由：}要求 $a>0$，$a\neq 1$，$N>0$，这里 $a=2$，$N=3$ 均满足。
		\item[\textbf{第二步（套用公式）：}]
			由恒等式 $a^{\log_a N}=N$，得 $2^{\log_2 3}=3$。
			\par\textbf{理由：}直接代入 $a=2$，$N=3$。
		\item[\textbf{第三步（给出答案）：}]
			故结果为 $3$。
			\par\textbf{理由：}恒等式成立，计算结束。
	\end{description}
	\textbf{\textcolor{CExample}{关键结论：}}
	$\boxed{2^{\log_2 3}=3}$。

	\medskip
	\textbf{\textcolor{CExample}{例2（稍变形）}}
	\textbf{\textcolor{CExample}{题目：}}
	计算 $9^{\log_3 5}$ 的值。
	\textbf{\textcolor{CExample}{解题步骤：}}
	\begin{description}[style=nextline, leftmargin=2.8em, labelsep=0.5em, itemsep=0.45em]
		\item[\textbf{第一步（底数变形）：}]
			将 $9$ 写成 $3^2$，则 $9^{\log_3 5} = (3^2)^{\log_3 5} = 3^{2\log_3 5}$。
			\par\textbf{理由：}利用幂的乘方运算法则：$(a^m)^n = a^{mn}$。
		\item[\textbf{第二步（指数处理）：}]
			利用对数性质，$2\log_3 5 = \log_3 (5^2) = \log_3 25$。
			\[
				3^{2\log_3 5} = 3^{\log_3 25}
			\]
			\par\textbf{理由：}对数运算法则：$n\log_a M = \log_a (M^n)$。
		\item[\textbf{第三步（套用恒等式）：}]
			由恒等式，$3^{\log_3 25} = 25$。
			\par\textbf{理由：}底数均为 $3$，真数 $25>0$，符合条件。
	\end{description}
	\textbf{\textcolor{CExample}{关键结论：}}
	$\boxed{9^{\log_3 5}=25}$。

	\medskip
	\textbf{\textcolor{CExample}{例3（综合应用）}}
	\textbf{\textcolor{CExample}{题目：}}
	解方程 $2^{\log_2 (x^2-3x+2)} = 4$。
	\textbf{\textcolor{CExample}{解题步骤：}}
	\begin{description}[style=nextline, leftmargin=2.8em, labelsep=0.5em, itemsep=0.45em]
		\item[\textbf{第一步（恒等化简）：}]
			左边指数底数与对数底数相同，若 $x^2-3x+2 > 0$，则恒等式成立，得 $2^{\log_2 (x^2-3x+2)} = x^2-3x+2$。
			\par\textbf{理由：}恒等式 $a^{\log_a N}=N$ 要求 $N>0$，此处 $N = x^2-3x+2$。
		\item[\textbf{第二步（转化方程）：}]
			原方程化为 $x^2-3x+2 = 4$，整理得 $x^2-3x-2=0$。
			\par\textbf{理由：}利用恒等式消去指数与对数。
		\item[\textbf{第三步（解二次方程）：}]
			解 $x^2-3x-2=0$，得 $x = \frac{3\pm\sqrt{17}}{2}$。
			\par\textbf{理由：}一元二次方程求根公式。
		\item[\textbf{第四步（检验定义域）：}]
			验证 $x$ 是否满足 $x^2-3x+2 > 0$。因为 $x^2-3x+2 = 4 >0$，所以两个解都满足条件。
			\par\textbf{理由：}原方程中的对数要求真数大于0，必须检验，避免增解。
	\end{description}
	\textbf{\textcolor{CExample}{关键结论：}}
	方程的解为 $\boxed{x = \frac{3\pm\sqrt{17}}{2}}$。
\end{examplebox}
